% =============================================================================
% INTERLUDIUM: DE FUGA
% Status: DRAFTED
% Placement: Between XVI and XVII (or around XVIII)
% Conventional narrative/thriller register — no typographic distortion.
% Creates contrast: over-stability after extreme fragmentation. Reads as
% almost uncanny in its coherence. Core ideas encoded in clean prose.
% =============================================================================

\chapter*{Interludium: De Fuga}
\addcontentsline{toc}{chapter}{Interludium. De Fuga}

\setstretch{1.03}

The door was marked before sunset.

Aisha saw it from across the narrow street, just as the last light slipped along the
plaster walls and gathered in the shallow cracks near the threshold. At first she
thought it was shadow, or some trick of the evening glow, but as she moved closer
the color resolved into something deliberate---thick, uneven, unmistakable.

Red.

Not painted. Applied.

She stopped a few paces away and did not touch it.

Behind her, the city continued as it always did at that hour. Vendors calling out the
last of their goods, shutters closing, the slow reshaping of noise into something quieter
but no less dense. Nothing in the rhythm of Córdoba suggested that anything had
changed.

Which meant everything had.

She turned the handle and entered.

Inside, the house was already in motion. Her father had not waited. Scrolls were
stacked in careful bundles across the low table, each tied with a strip of linen and
marked in a hand that was more compressed than usual, as if time itself had narrowed
the space available for writing.

Averroes stood near the center of the room, holding a manuscript open but not
reading it.

``You saw it,'' he said, without turning.

Aisha closed the door behind her.

``Yes.''

``They chose visibility,'' she added. ``Not force.''

Now he turned.

``That is force,'' he said. ``Applied differently.''

From the courtyard came the faint sound of movement---someone passing too slowly,
or stopping where there was no reason to stop. Aisha moved to the window and
shifted the lattice just enough to see.

Two men stood across the street, speaking in low tones. They were not dressed as
officials. That was part of the point.

``They're not hiding it,'' she said.

``No,'' Averroes replied. ``They don't need to.''

A sound at the door. Not a knock. Something lighter. A signal, perhaps.

Aisha moved before he did, crossing the room and opening it just enough to admit
a figure from the street.

Zaynab slipped inside without a word, her expression composed but tight at the edges.

``They've started,'' she said.

``Where?'' Aisha asked.

``The market first. Then the western quarter.'' Zaynab set a small satchel on the
table. ``Confiscations. Not systematic yet. But moving.''

Averroes closed the manuscript in his hands.

``Then we leave tonight.''

She gathered the bundles. Their weight distributed between them.

``Father,'' she said.

He looked at her.

``If they had not done this---would you have stayed?''

The question hung between them, more dangerous than the noise outside.

``For a time,'' he said. ``Long enough to finish what could be finished.''

``And now?''

``Now,'' he said, ``it will be finished elsewhere.''

Zaynab moved to the door, listening.

``Wait,'' she said.

They did.

Footsteps passed. Slower this time. Deliberate.

Then silence.

``Now,'' Zaynab said.

They stepped into the street together.

The red mark on the door caught the last of the fading light, darker now, almost
indistinguishable from shadow.

Aisha did not look back.

They moved north.

Behind them, Córdoba continued---its streets filling, emptying, reforming---unchanged
in appearance, altered in structure.

Ahead, the road narrowed, then opened, then disappeared into darkness.

They followed it anyway.
